%% kapa: exam=2 round=20 year=2009 period=1 subject=appraisal_practice subject_ko=감정평가실무 seq=1 number=1 points=40
%% source: Q-Net 감정평가사 기출문제 (cst00309) artlSeq=5024180 file=extracted/1008/5024180/2009년 제20회 감정평가사 2차시험 기출문제/1교시.hwp
%% conversion: kordoc 4.24.51 -> tools/gen2.py (auto); review=proofread-sonnet; flags=none
\begin{longquestion}{【문제 1】}
자동차부품업체를 운영하고 있는 김갑동 사장은 공장을 증설하기 위하여 임야를 매입하고, 자금마련을 위해 개발단계별로 담보대출을 신청하려 한다. 주어진 조건과 자료를 참고하여 다음 물음에 답하시오. (40점) \par
\medskip\noindent (1) 2009.01.01을 가격시점으로 하여 토지를 평가하시오. (5점) \par
\medskip\noindent (2) 2009.03.31을 가격시점으로 하여 토지를 평가하시오. (10점) \par
\medskip\noindent (3) 2009.09.06을 가격시점으로 하여 공장을 평가하시오. (25점) \par
<자료 1> 평가의뢰 내역 \par
1) 평가물건 \par
① 토지: C시 Y읍 S리 산11번지 중 김갑동 소유지분 \par
② 건물: 위 지상 소재 건물 \par
2) 평가목적: 담보 \par
<자료 2> 2009.01.01 가격시점 관련 사항 \par
1) 사전조사사항 \par
① 토지등기부등본 \par
{\small\centering
\begin{tabular}{|c|c|c|c|}
\hline
\makecell{소재지번} & \makecell{지목} & \makecell{면적} & \makecell{소유자} \\
\hline
\makecell{C시 Y읍 S리 산11} & \makecell{임야} & \makecell{23,955m\textsuperscript{2}} & \makecell{공유자 지분 3분의 1 김갑동\newline{}지분 3분의 1 이갑동\newline{}지분 3분의 1 박갑동} \\
\hline
\end{tabular}\par}

② 건물등기부등본 및 건축물대장등본: 해당사항 없음 \par
③ 토지대장등본: 등기부와 동일 \par
④ 지적도: 지적분할신청 중으로 발급받지 못함 \par
⑤ 토지이용계획확인원: 계획관리지역, 준보전산지 \par
⑥ 공장신설승인신청서 사본(요약) \par
{\small\centering
\begin{tabular}{|c|c|c|c|c|}
\hline
\makecell{소재지번} & \makecell{용도지역} & \makecell{공장용지\newline{}면적} & \makecell{제조시설\newline{}면적} & \makecell{부대시설\newline{}면적} \\
\hline
\makecell{C시 Y읍 S리\newline{}산11번지\newline{}(분할 후 11번지)} & \makecell{계획관리\newline{}지역} & \makecell{7,780m\textsuperscript{2}} & \makecell{2,000m\textsuperscript{2}} & \makecell{500m\textsuperscript{2}} \\
\hline
\end{tabular}\par}

\textbf{※ 분할 후 11-3번지는 진입도로로 조성할 것이며 토지 가분할 측량성과도와 같이 분할예정임.} \par
⑦ 토지 가분할 측량성과도 \par
\qfigure{figures/01-1.png}
2) 현장조사사항 \par
① 대상토지인 산11번지는 왕복 2차선 국도변에 위치한 남서향 완경사의 임야로 대부분 활잡목인 임지상의 임목은 별도의 평가가치는 없는 것으로 판단되었으며 부근은 국도주변 중소규모 공장 및 농경지대임. \par
② 대상토지는 토지분할 및 공장신설승인신청서가 곧 수리될 예정인 것으로 탐문되었음. \par
3) 심사평가사의 심사의견 \par
당해지역은 2009.01.01 기준으로 관리지역 세분화가 시행되었고 임야의 경우 산지번에서 등록전환, 분할측량 등의 경우에는 면적이 달라질 수 있음. \par
<자료 3> 2009.03.31 가격시점 관련 사항 \par
1) 사전조사사항 \par
① 토지등기부등본: C시 Y읍 S리 산11에서 C시 Y읍 S리 11로 등록전환되고 면적은 23,940m\textsuperscript{2}로 변경되었으며, 지목과 소유자는 동일함. \par
② 건물등기부등본 및 건축물대장등본: 해당사항 없음. \par
③ 토지대장등본 \par
{\small\centering
\resizebox{\ifdim\width>\linewidth \linewidth\else\width\fi}{!}{%
\begin{tabular}{|c|c|c|c|c|c|}
\hline
\multirow{2}{*}{\makecell{토지소재}} & \multirow{2}{*}{\makecell{지번}} & \multicolumn{3}{|c|}{\makecell{토지표시}} & \multirow{2}{*}{\makecell{소유자}} \\
\cline{3-5}
 &  & \makecell{지목} & \makecell{면적} & \makecell{사유} &  \\
\hline
\makecell{C시 Y읍 S리} & \makecell{11} & \makecell{임야} & \makecell{23,940m\textsuperscript{2}} & \makecell{2009년 03월 01일\newline{}산11에서 등록전환} & \makecell{공유자\newline{}지분3분의1 김갑동\newline{}지분3분의1 이갑동\newline{}지분3분의1 박갑동} \\
\hline
\makecell{C시 Y읍 S리} & \makecell{11} & \makecell{임야} & \makecell{7,780m\textsuperscript{2}} & \makecell{2009년 03월 02일\newline{}분할되어 본번에 -1, -2, -3을 부함} & \makecell{김갑동} \\
\hline
\end{tabular}}\par}

{\small\centering
\resizebox{\ifdim\width>\linewidth \linewidth\else\width\fi}{!}{%
\begin{tabular}{|c|c|c|c|c|c|}
\hline
\multirow{2}{*}{\makecell{토지소재}} & \multirow{2}{*}{\makecell{지번}} & \multicolumn{3}{|c|}{\makecell{토지표시}} & \multirow{2}{*}{\makecell{소유자}} \\
\cline{3-5}
 &  & \makecell{지목} & \makecell{면적} & \makecell{사유} &  \\
\hline
\makecell{C시 Y읍 S리} & \makecell{11-1} & \makecell{임야} & \makecell{7,780m\textsuperscript{2}} & \makecell{2009년 03월 02일\newline{}11번지에서 분할} & \makecell{이갑동} \\
\hline
\end{tabular}}\par}

{\small\centering
\resizebox{\ifdim\width>\linewidth \linewidth\else\width\fi}{!}{%
\begin{tabular}{|c|c|c|c|c|c|}
\hline
\multirow{2}{*}{\makecell{토지소재}} & \multirow{2}{*}{\makecell{지번}} & \multicolumn{3}{|c|}{\makecell{토지표시}} & \multirow{2}{*}{\makecell{소유자}} \\
\cline{3-5}
 &  & \makecell{지목} & \makecell{면적} & \makecell{사유} &  \\
\hline
\makecell{C시 Y읍 S리} & \makecell{11-2} & \makecell{임야} & \makecell{7,780m\textsuperscript{2}} & \makecell{2009년 03월 02일\newline{}11번지에서 분할} & \makecell{박갑동} \\
\hline
\end{tabular}}\par}

{\small\centering
\resizebox{\ifdim\width>\linewidth \linewidth\else\width\fi}{!}{%
\begin{tabular}{|c|c|c|c|c|c|}
\hline
\multirow{2}{*}{\makecell{토지소재}} & \multirow{2}{*}{\makecell{지번}} & \multicolumn{3}{|c|}{\makecell{토지표시}} & \multirow{2}{*}{\makecell{소유자}} \\
\cline{3-5}
 &  & \makecell{지목} & \makecell{면적} & \makecell{사유} &  \\
\hline
\makecell{C시 Y읍 S리} & \makecell{11-3} & \makecell{임야} & \makecell{600m\textsuperscript{2}} & \makecell{2009년 03월 02일\newline{}11번지에서 분할} & \makecell{공유자\newline{}지분3분의1 김갑동\newline{}지분3분의1 이갑동\newline{}지분3분의1 박갑동} \\
\hline
\end{tabular}}\par}

④ 토지이용계획확인원 \par
{\small\centering
\begin{tabular}{|c|c|c|c|c|}
\hline
\makecell{토지소재} & \makecell{지번} & \makecell{지목} & \makecell{면적} & \makecell{토지이용계획사항} \\
\hline
\makecell{C시 Y읍 S리} & \makecell{11} & \makecell{임야} & \makecell{7,780m\textsuperscript{2}} & \makecell{계획관리지역, 준보전산지} \\
\hline
\makecell{C시 Y읍 S리} & \makecell{11-1} & \makecell{임야} & \makecell{7,780m\textsuperscript{2}} & \makecell{계획관리지역, 준보전산지} \\
\hline
\makecell{C시 Y읍 S리} & \makecell{11-2} & \makecell{임야} & \makecell{7,780m\textsuperscript{2}} & \makecell{계획관리지역, 준보전산지} \\
\hline
\makecell{C시 Y읍 S리} & \makecell{11-3} & \makecell{임야} & \makecell{600m\textsuperscript{2}} & \makecell{계획관리지역, 준보전산지} \\
\hline
\end{tabular}\par}

⑤ 지적도 및 기타사항: 토지 가분할 측량성과도와 같이 분할되어 확정되었으며, 공장신설건은 2009.03.10자로 신청서와 같이 승인되었음. \par
2) 현장조사사항 \par
① 대상토지는 인접토지와 평탄하게 공장부지 조성공사(조경․바닥포장공사는 착수하지 않았음) 및 접면도로 포장공사가 완료되어 있었음. \par
② 현장조사시 제시받은 공장부지 조성원가 자료는 아래와 같음. \par
{\small\centering
\begin{tabular}{|c|c|}
\hline
\makecell{구 분} & \makecell{금 액(단위: 원)} \\
\hline
\makecell{가설 및 토공사} & \makecell{45,000,000} \\
\hline
\makecell{자재 및 운반비} & \makecell{150,000,000} \\
\hline
\makecell{옹벽공사} & \makecell{30,000,000} \\
\hline
\makecell{조경․바닥포장공사} & \makecell{55,000,000} \\
\hline
\makecell{기타 제간접경비 등} & \makecell{72,000,000} \\
\hline
\end{tabular}\par}

\textbf{※ 접면도로 포장비는 포함되어 있지 않고 별도 고려하지 아니함.} \par
③ 대상토지의 공장용지부분에 건물신축을 위해 임시사용승인을 받은 경량철골조 철판지붕 단층 작업장(바닥면적: 100m\textsuperscript{2})이 소재함. \par
3) 심사평가사의 심사의견 \par
① 제시한 조성공사비의 대부분은 적정하나 자재비가 일시 폭등한 시점에 공사가 이뤄져 자재 및 운반비는 통상적인 공사에 비해 50\% 정도 높은 것으로 보이니 가격검토시 이를 고려할 것 (단, 기타 간접제경비 등은 제시금액으로 할 것) \par
② 만약 제시외 건물의 토지에 대한 영향정도를 파악할 경우 건물의 바닥면적 만큼만 고려할 것 \par
③ 막다른 길이 있는 각지의 도로접면은 한면으로 인식할 것 \par
<자료 4> 2009.09.06 가격시점 관련 사항 \par
1) 사전조사사항 \par
① 토지등기부등본 및 토지대장등본: 2009.03.31 토지대장과 동일 \par
② 건물등기부등본: 해당사항 없음 \par
③ 토지이용계획확인원 및 지적도: 종전과 동일 \par
④ 건축물대장등본: 소유자는 김갑동임 \par
{\small\centering
\begin{tabular}{|c|c|c|c|c|}
\hline
\makecell{대지위치} & \makecell{지번} & \makecell{대지면적} & \makecell{건축면적} & \makecell{사용승인일자} \\
\hline
\makecell{C시 Y읍 S리} & \makecell{11} & \makecell{7,780m\textsuperscript{2}} & \makecell{2,500m\textsuperscript{2}} & \makecell{2009.09.06} \\
\hline
\makecell{구 분} & \makecell{층별} & \makecell{구조} & \makecell{용도} & \makecell{면적} \\
\hline
\makecell{주1} & \makecell{1} & \makecell{일반철골구조} & \makecell{공장} & \makecell{2,000m\textsuperscript{2}} \\
\hline
\makecell{부1} & \makecell{1} & \makecell{일반철골구조} & \makecell{사무실} & \makecell{500m\textsuperscript{2}} \\
\hline
\end{tabular}\par}

\textbf{※ 5일 이내에 토지지목 변경을 조건부로 한 사용승인이었음.} \par
⑤ 기계기구 의뢰목록 \par
{\small\centering
\begin{tabular}{|c|c|c|c|c|}
\hline
\makecell{구분} & \makecell{기계명} & \makecell{수량} & \makecell{제작 및 구입일자} & \makecell{구입가격(원/대)} \\
\hline
\makecell{1} & \makecell{CNC M/C\newline{}(수치제어선반)} & \makecell{2} & \makecell{수입신고서 참조} & \makecell{수입신고서 참조} \\
\hline
\makecell{2} & \makecell{선 반} & \makecell{3} & \makecell{2009.01.01} & \makecell{50,000,000} \\
\hline
\makecell{3} & \makecell{Air Compressor\newline{}(컴프레셔)} & \makecell{1} & \makecell{2008.08.01} & \makecell{12,000,000} \\
\hline
\end{tabular}\par}

2) 현장조사사항 \par
① 공장부지 조성공사는 완료되어 있었음. \par
② 현장조사시 제시받은 건물공사비 내역서 자료는 다음과 같음. \par
{\small\centering
\begin{tabular}{|c|c|c|}
\hline
\makecell{구분} & \makecell{공장동(단위: 원)} & \makecell{사무실동(단위: 원)} \\
\hline
\makecell{기초공사} & \makecell{30,000,000} & \makecell{5,000,000} \\
\hline
\makecell{옹벽공사} & \makecell{30,000,000} & \makecell{--} \\
\hline
\makecell{철골 및 철근콘크리트 공사} & \makecell{250,000,000} & \makecell{30,000,000} \\
\hline
\makecell{조적 및 벽체공사} & \makecell{120,000,000} & \makecell{15,000,000} \\
\hline
\makecell{창호 및 지붕공사} & \makecell{100,000,000} & \makecell{13,000,000} \\
\hline
\makecell{미장, 타일, 도장, 위생 및 냉난방공사 등} & \makecell{170,000,000} & \makecell{31,000,000} \\
\hline
\makecell{일반관리비 등(간접경비)} & \makecell{50,000,000} & \makecell{25,000,000} \\
\hline
\makecell{설계, 감리, 전기기본공사} & \makecell{100,000,000} & \makecell{19,000,000} \\
\hline
\makecell{수배전설비(100kw)} & \makecell{150,000,000} & \makecell{--} \\
\hline
\makecell{크레인설비(20ton)} & \makecell{15,000,000} & \makecell{--} \\
\hline
\end{tabular}\par}

③ 기계장치는 신규 설치되어 정상가동되고 있으며 도입기계 중 CNC M/C 1대는 향후 증설을 예상하여 도입하였으나 설치하지 않고 보관 중으로 증설시기는 미정이며, 목록에 포함되지 않은 기계기구의 제작 및 구입일자는 선반과 동일한 것으로 조사되었음. \par
④ 도입기계 관련 수입신고서(요약) \par
{\small\centering
\begin{tabular}{|c|c|c|c|}
\hline
\makecell{신고일} & \makecell{입항일} & \makecell{반입일} & \makecell{적출국} \\
\hline
\makecell{2009.02.01} & \makecell{2009.01.05} & \makecell{2009.01.08} & \makecell{JP(JAPAN)} \\
\hline
\makecell{품명} & \makecell{수량} & \makecell{단가(USD)} & \makecell{금액(USD)} \\
\hline
\makecell{CNC M/C} & \makecell{2U} & \makecell{100,000} & \makecell{200,000} \\
\hline
\multirow{2}{*}{\makecell{과세가격(CIF)}} & \makecell{\$200,000} & \multirow{2}{*}{\makecell{원산지 표시}} & \multirow{2}{*}{\makecell{JP-Y-Z-N}} \\
\cline{2-2}
 & \makecell{￦280,304,000} &  &  \\
\hline
\makecell{세종\newline{}관\newline{}농\newline{}부} & \makecell{세율\newline{}8.00\newline{}20.00\newline{}10.00} & \makecell{감면율\newline{}50.000} & \makecell{세액\newline{}11,212,160\newline{}2,242,432\newline{}29,375,859} \\
\hline
\makecell{결재금액} & \makecell{CIF-USD200,000} & \makecell{환율} & \makecell{1,401.52} \\
\hline
\end{tabular}\par}

3) 심사평가사의 심사의견 \par
업자가 제시한 기계기구 구입가격 및 건물공사비 내역서의 금액은 적정한 것으로 보이나 일부 항목은 건물공사비 산입의 적정성을 재검토하고 특히 기계기구 의뢰목록은 재작성해야 할 것이라는 의견을 제시함. \par
<자료 5> 가격결정을 위한 참고자료 \par
1) 표준지 공시지가 현황(현장조사일과 감정평가서 작성완료일은 동일하고 공시지가 공시기준일은 매년 1월 1일, 공시일은 매년 3월 1일임) \par
{\small\centering
\resizebox{\ifdim\width>\linewidth \linewidth\else\width\fi}{!}{%
\begin{tabular}{|c|c|c|c|c|c|c|c|c|c|c|}
\hline
\multirow{2}{*}{\makecell{일련번호}} & \multirow{2}{*}{\makecell{소재지}} & \multirow{2}{*}{\makecell{면적\newline{}(m\textsuperscript{2})}} & \multirow{2}{*}{\makecell{지목}} & \multirow{2}{*}{\makecell{이용\newline{}상황}} & \multicolumn{2}{|c|}{\makecell{용도\newline{}지역}} & \multirow{2}{*}{\makecell{도로\newline{}교통}} & \multirow{2}{*}{\makecell{형상 및\newline{}지세}} & \multicolumn{2}{|c|}{\makecell{공시지가(원/m\textsuperscript{2})}} \\
\cline{6-7} \cline{10-11}
 &  &  &  &  & \makecell{2008년도} & \makecell{2009\newline{}년도} &  &  & \makecell{2008\newline{}년도} & \makecell{2009\newline{}년도} \\
\hline
\makecell{1} & \makecell{C시 Y읍 S리 산20} & \makecell{17,345} & \makecell{임야} & \makecell{임야} & \makecell{관리} & \makecell{계획\newline{}관리} & \makecell{맹지} & \makecell{부정형\newline{}완경사} & \makecell{51,000} & \makecell{50,000} \\
\hline
\makecell{2} & \makecell{C시 Y읍 J리 산17} & \makecell{22,915} & \makecell{임야} & \makecell{임야} & \makecell{관리} & \makecell{보전관리} & \makecell{세로가} & \makecell{부정형\newline{}완경사} & \makecell{39,000} & \makecell{38,000} \\
\hline
\makecell{3} & \makecell{C시 Y읍 S리 107} & \makecell{8,950} & \makecell{공장\newline{}용지} & \makecell{공업용} & \makecell{관리} & \makecell{계획\newline{}관리} & \makecell{소로\newline{}한면} & \makecell{부정형\newline{}평지} & \makecell{151,000} & \makecell{150,000} \\
\hline
\makecell{4} & \makecell{C시 Y읍 S리 55} & \makecell{2,235} & \makecell{잡종지} & \makecell{상업용} & \makecell{관리} & \makecell{계획\newline{}관리} & \makecell{소로\newline{}한면} & \makecell{장방형\newline{}평지} & \makecell{223,000} & \makecell{220,000} \\
\hline
\end{tabular}}\par}

2) 적용할 지가변동률(월말에 해당 월 변동률을 발표한다고 간주, 단위: \%) \par
① 2008년도 \par
{\small\centering
\begin{tabular}{|c|c|c|c|c|c|}
\hline
\makecell{구 분} & \makecell{공업지역} & \makecell{관리지역} & \makecell{농림지역} & \makecell{임야} & \makecell{공업용} \\
\hline
\makecell{2008.01.01～12.31} & \makecell{-1.179} & \makecell{-1.245} & \makecell{-1.377} & \makecell{-1.154} & \makecell{-0.912} \\
\hline
\makecell{2008.12.01～12.31} & \makecell{-0.179} & \makecell{-0.389} & \makecell{-0.247} & \makecell{-0.169} & \makecell{-0.088} \\
\hline
\end{tabular}\par}

② 2009년도 \par
{\small\centering
\begin{tabular}{|c|c|c|c|c|c|}
\hline
\makecell{구 분} & \makecell{공업지역} & \makecell{보전관리} & \makecell{계획관리} & \makecell{임야} & \makecell{공업용} \\
\hline
\makecell{2009.01.01～03.31} & \makecell{-0.697} & \makecell{-0.765} & \makecell{-0.454} & \makecell{-0.667} & \makecell{-0.546} \\
\hline
\makecell{2009.01.01～09.06} & \makecell{0.998} & \makecell{0.996} & \makecell{0.997} & \makecell{0.988} & \makecell{0.917} \\
\hline
\end{tabular}\par}

3) 지역요인: 동일함. \par
4) 개별요인: 이용상황이 동일하면 별도의 지목감가는 하지 아니함. \par
① 도로접면 \par
{\small\centering
\begin{tabular}{|c|c|c|c|c|}
\hline
\makecell{구분} & \makecell{소로한면} & \makecell{세로가} & \makecell{세로(불)} & \makecell{맹지} \\
\hline
\makecell{소로한면} & \makecell{1.00} & \makecell{0.93} & \makecell{0.86} & \makecell{0.83} \\
\hline
\makecell{세로가} & \makecell{1.07} & \makecell{1.00} & \makecell{0.92} & \makecell{0.89} \\
\hline
\makecell{세로(불)} & \makecell{1.16} & \makecell{1.09} & \makecell{1.00} & \makecell{0.96} \\
\hline
\makecell{맹지} & \makecell{1.20} & \makecell{1.12} & \makecell{1.04} & \makecell{1.00} \\
\hline
\end{tabular}\par}

② 형상 \par
{\small\centering
\begin{tabular}{|c|c|c|c|c|}
\hline
\makecell{구분} & \makecell{정방형} & \makecell{장방형} & \makecell{사다리형} & \makecell{부정형} \\
\hline
\makecell{정방형} & \makecell{1.00} & \makecell{0.99} & \makecell{0.98} & \makecell{0.95} \\
\hline
\makecell{장방형} & \makecell{1.01} & \makecell{1.00} & \makecell{0.99} & \makecell{0.96} \\
\hline
\makecell{사다리형} & \makecell{1.02} & \makecell{1.01} & \makecell{1.00} & \makecell{0.97} \\
\hline
\makecell{부정형} & \makecell{1.05} & \makecell{1.04} & \makecell{1.03} & \makecell{1.00} \\
\hline
\end{tabular}\par}

③ 지세 \par
{\small\centering
\begin{tabular}{|c|c|c|}
\hline
\makecell{구분} & \makecell{평지} & \makecell{완경사} \\
\hline
\makecell{평지} & \makecell{1.00} & \makecell{0.97} \\
\hline
\makecell{완경사} & \makecell{1.03} & \makecell{1.00} \\
\hline
\end{tabular}\par}

④ 2009.03.31 기준 C시 Y읍 S리 11번지 토지의 성숙도 비교치 \par
{\small\centering
\begin{tabular}{|c|c|c|c|c|}
\hline
\makecell{대상토지} & \makecell{표준지 3} & \makecell{표준지 4} & \makecell{거래사례} & \makecell{평가선례} \\
\hline
\makecell{1.00} & \makecell{1.10} & \makecell{1.10} & \makecell{0.50} & \makecell{0.90} \\
\hline
\end{tabular}\par}

5) 거래사례 및 평가선례 \par
① 거래사례 \par
{\small\centering
\begin{tabular}{|c|c|c|c|c|c|c|c|c|}
\hline
\makecell{소재지} & \makecell{지목} & \makecell{면적\newline{}(m\textsuperscript{2})} & \makecell{이용\newline{}상황} & \makecell{용도지역} & \makecell{도로교통} & \makecell{형상 및\newline{}지세} & \makecell{단가\newline{}(원/m\textsuperscript{2})} & \makecell{거래시점} \\
\hline
\makecell{C시 Y읍\newline{}S리 산11} & \makecell{임야} & \makecell{7,985} & \makecell{임야} & \makecell{관리} & \makecell{소로한면} & \makecell{부정형완경사} & \makecell{110,000} & \makecell{2008.12.01} \\
\hline
\end{tabular}\par}

② 평가선례: 유사사례가 많으나 대표적인 것만 제시함. \par
{\small\centering
\begin{tabular}{|c|c|c|c|c|c|c|c|c|}
\hline
\makecell{소재지} & \makecell{지목} & \makecell{면적\newline{}(m\textsuperscript{2})} & \makecell{이용\newline{}상황} & \makecell{용도지역} & \makecell{도로교통} & \makecell{형상 및\newline{}지세} & \makecell{단가\newline{}(원/m\textsuperscript{2})} & \makecell{가격시점} \\
\hline
\makecell{C시 Y읍\newline{}S리 산22} & \makecell{임야} & \makecell{7,890} & \makecell{공장\newline{}예정지} & \makecell{계획관리} & \makecell{소로한면} & \makecell{부정형완경사} & \makecell{120,000} & \makecell{2009.01.01} \\
\hline
\end{tabular}\par}

③ 심사평가사의 심사의견 \par
수집한 자료들 중 평가선례는 적정하나 거래사례는 개발이익의 상당부분이 매도자에게 귀속된 것으로 보이고 공장예정지인 평가선례는 개별요인에서 성숙도를 보정해야 한다는 의견을 제시함. \par
6) 원가법에 의한 평가시 투하자금에 대한 기간이자는 고려하지 아니함. \par
7) 건물평가자료: 제시자료를 활용하되, 내용연수는 35년을 적용할 것 \par
8) 기계기구 평가자료 \par
① 내용연수는 15년, 최종 잔가율은 10\%를 적용 \par
② 도입기계 관련 자료 \par
○ CIF, 원산지 화폐를 기준하고 국내 시장가격은 고려하지 아니함. \par
○ 기계가격보정지수: 1.0 \par
○ 외화환산율 \par
{\small\centering
\begin{tabular}{|c|c|c|c|c|}
\hline
\makecell{적용시점} & \makecell{통화} & \makecell{해당통화당\newline{}미(달러)} & \makecell{미(달러)당\newline{}해당통화} & \makecell{해당통화당\newline{}한국(원)} \\
\hline
\makecell{2009. 01} & \makecell{JPY} & \makecell{0.7150(100엔당)} & \makecell{139.8601} & \makecell{1,409.10(100엔당)} \\
\hline
\makecell{2009. 02} & \makecell{JPY} & \makecell{0.7532(100엔당)} & \makecell{132.7669} & \makecell{1,425.05(100엔당)} \\
\hline
\makecell{2009. 08} & \makecell{JPY} & \makecell{0.7635(100엔당)} & \makecell{130.9758} & \makecell{1,405.22(100엔당)} \\
\hline
\end{tabular}\par}

○ 도입부대비 \par
설치비는 도입가격의 1.5\%, L/C개설비 등 기타 부대비용은 도입가격의 3\%를 적용하고 세율, 감면율 등은 도입시점과 동일하게 적용 \par
○ 정률법에 의한 잔존가치율(내용연수는 15년, 최종 잔가율은 10\%) \par
{\small\centering
\begin{tabular}{|c|c|c|c|c|}
\hline
\makecell{경과년수} & \makecell{1} & \makecell{2} & \makecell{3} & \makecell{4} \\
\hline
\makecell{잔존가치율} & \makecell{0.858} & \makecell{0.736} & \makecell{0.631} & \makecell{0.541} \\
\hline
\end{tabular}\par}

\end{longquestion}
